\documentclass{article}

\usepackage{amsmath,amssymb}
\usepackage{xurl}
\usepackage[hidelinks]{hyperref}
\setlength{\emergencystretch}{2em}

\input{command}

\title{Normalization of the tensor printed in CPSV16 Example 4.12}
\author{}
\date{2026-08-10}

\begin{document}
\maketitle
\gapnote{false-source}{open}

\section{Printed representative}

Example~4.12 of Ref.~\cite{Cirac2016MPDO_arXiv}
(\path{Papers/1606.00608/MPDO-22-12-17-2.tex:932--939}) prints a tensor $M$
with physical dimension $d=2$, virtual dimension $D=2$, and components
\begin{align}
    1
        &= M_{00}^{00}
         = M_{00}^{11}
         = M_{11}^{00}
         = -M_{11}^{11},
    \label{eq:cpsv412_printed_components}
\end{align}
with all other components equal to zero. In matrix notation,
\begin{align}
    M^{00}
        &= \Id,
    \label{eq:cpsv412_printed_matrix_00}\\
    M^{11}
        &= \sigma_z,
    \label{eq:cpsv412_printed_matrix_11}\\
    M^{01}
        &= M^{10}
         = 0,
    \label{eq:cpsv412_printed_off_diagonal_matrices}
\end{align}
where
\begin{align}
    \sigma_z
        &= \ketbra{0}{0}-\ketbra{1}{1}.
    \label{eq:cpsv412_pauli_z}
\end{align}
Let $\rho^{(N)}(M)$ denote the periodic length-$N$ operator obtained by closing
the virtual indices of $M$. For every $N>0$, the printed tensor gives
\begin{align}
    \rho^{(N)}(M)
        &= \Id^{\otimes N}+\sigma_z^{\otimes N}.
    \label{eq:cpsv412_literal_periodic_family}
\end{align}
Its trace is
\begin{align}
    \tr\!\left(\rho^{(N)}(M)\right)
        &= 2^N,
    \label{eq:cpsv412_literal_trace}
\end{align}
not $1$. Thus the printed $M$ is an unnormalized tensor representative,
although division of the resulting positive operator by $2^N$ gives a
normalized density operator.

\section{The fixed-representative mismatch}

Define the physical-trace transfer matrix of the printed tensor by
$\mathcal T_M=M^{00}+M^{11}$. Equations~\ref{eq:cpsv412_printed_matrix_00}
and \ref{eq:cpsv412_printed_matrix_11} give
\begin{align}
    \mathcal T_M
        &= M^{00}+M^{11}
         = \Id+\sigma_z
         = \begin{pmatrix}
               2 & 0\\
               0 & 0
           \end{pmatrix}.
    \label{eq:cpsv412_literal_transfer}
\end{align}
Consequently,
\begin{align}
    \mathcal T_M^2
        &= 2\mathcal T_M,
    \label{eq:cpsv412_literal_transfer_scaled_idempotence}\\
    \mathcal T_M^2
        &\neq \mathcal T_M.
    \label{eq:cpsv412_literal_transfer_not_idempotent}
\end{align}
Equation~\ref{eq:cpsv412_literal_transfer_scaled_idempotence} proves the
project's scale-invariant relation with scale $2$, but
Eq.~\ref{eq:cpsv412_literal_transfer_not_idempotent} contradicts the literal
zero-correlation-length equation in Definition~4.2 of
Ref.~\cite{Cirac2016MPDO_arXiv}. The trace-preserving maps $\mathcal S$ and
$\mathcal T$ in Definition~4.1 of Ref.~\cite{Cirac2016MPDO_arXiv} also force
equality of the one-site and two-site traces for every virtual boundary.
Equivalently, they require literal idempotence of $\mathcal T_M$. The printed
representative therefore cannot satisfy that fixed-tensor renormalization
condition.

This failure does not contradict the normalized density family. Define the
rescaled representative by
\begin{align}
    \widehat M
        &= \frac{1}{2}M.
    \label{eq:cpsv412_normalized_representative}
\end{align}
Its transfer matrix satisfies
\begin{align}
    \mathcal T_{\widehat M}
        &= \frac{1}{2}\mathcal T_M,
    \label{eq:cpsv412_normalized_transfer}\\
    \mathcal T_{\widehat M}^2
        &= \mathcal T_{\widehat M}.
    \label{eq:cpsv412_normalized_transfer_idempotent}
\end{align}
The periodic family of $\widehat M$ is
\begin{align}
    \rho^{(N)}(\widehat M)
        &= 2^{-N}\left(\Id^{\otimes N}+\sigma_z^{\otimes N}\right).
    \label{eq:cpsv412_normalized_periodic_family}
\end{align}
This is precisely the normalized density family associated with the printed
tensor. The representative $\widehat M=(1/2)M$ is therefore the likely tensor
intended by the assertion at source line~938 that trace-preserving maps exist
\cite{Cirac2016MPDO_arXiv}. The paper does not include this factor in the
printed tensor components or explicitly change representatives before making
that assertion.

\section{Formalization boundary}

The current formalization keeps the two representatives separate. For the
literal printed tensor $M$, it proves the exact positive-length formula,
positivity, one-site loss of the $\sigma_z$ term, saturation of the area law,
zero correlation length only in the project's scale-invariant sense, explicit
failure of literal zero correlation length, and failure of the
fixed-representative trace-preserving renormalization condition.

For $\widehat M=(1/2)M$, define the parity function
$\pi(a,b)=a\mathbin{\oplus}b$ and the Kraus operators
\begin{align}
    C_{a,b}
        &= \ketbra{\pi(a,b)}{a,b},
    \label{eq:cpsv412_coarse_graining_kraus}\\
    R_{a,b}
        &= 2^{-1/2}\ketbra{a,b}{\pi(a,b)}.
    \label{eq:cpsv412_refinement_kraus}
\end{align}
Each parity fibre has two elements. Therefore,
\begin{align}
    \sum_{a,b} C_{a,b}^{\dagger}C_{a,b}
        &= \Id_4,
    \label{eq:cpsv412_coarse_graining_trace_preserving}\\
    \sum_{a,b} R_{a,b}^{\dagger}R_{a,b}
        &= \Id_2.
    \label{eq:cpsv412_refinement_trace_preserving}
\end{align}
The corresponding completely positive maps are trace preserving. The local
tensor products satisfy
\begin{align}
    \widehat M^{aa}\widehat M^{bb}
        &= \frac{1}{2}
           \widehat M^{\pi(a,b)\pi(a,b)},
    \label{eq:cpsv412_normalized_tensor_parity_product}
\end{align}
and the off-diagonal physical letters vanish. Direct substitution therefore
proves both channel equations of Definition~4.1 of
Ref.~\cite{Cirac2016MPDO_arXiv} for every virtual boundary matrix. These
results are recorded by
\leanid{MPOTensor.CPSVExample412NormalizedRFP.Mhat},
\leanid{MPOTensor.CPSVExample412NormalizedRFP.mpo_Mhat_eq_normalizedMPO_M},
\leanid{MPOTensor.CPSVExample412NormalizedRFP.Mhat_isMPDO}, and
\leanid{MPOTensor.CPSVExample412NormalizedRFP.Mhat_isRFPViaTS}.

This conclusion concerns only the two channel equations. It does not resolve a
second normalization tension in the source. As a doubled-index matrix product
state, $M$ has two bond-one sectors with physical vectors
$(1,0,0,1)$ and $(1,0,0,-1)$. Division by $\sqrt{2}$ gives their normal
representatives. Consequently, the two horizontal normal-block weights of
$\widehat M$ are both $1/\sqrt{2}$. Neither has modulus one, whereas source
line~246 requires at least one unit-modulus weight
\cite{Cirac2016MPDO_arXiv}. The proved bare channel predicate must therefore
not be presented as a proof that $\widehat M$ satisfies every standing
canonical-form convention in Definition~4.1 of
Ref.~\cite{Cirac2016MPDO_arXiv}. This boundary is also recorded in
\path{docs/paper-gaps/cpsv16_unit_weight_rfp_scale_tension.tex}.

The independent four-site entropy calculation is also complete. For the
normalized state, regroup the sites as $A=\{0\}$, $B=\{1,3\}$, and
$C=\{2\}$. The $AB$ and $BC$ marginals are $\Id_8/8$, the $B$ marginal is
$\Id_4/4$, and the full state is uniform on eight even-parity configurations.
Thus,
\begin{align}
    S(ABC)
        &= 3\log 2,
    \label{eq:cpsv412_entropy_abc}\\
    S(B)
        &= 2\log 2,
    \label{eq:cpsv412_entropy_b}\\
    S(AB)
        &= 3\log 2,
    \label{eq:cpsv412_entropy_ab}\\
    S(BC)
        &= 3\log 2.
    \label{eq:cpsv412_entropy_bc}
\end{align}
The strong-subadditivity defect is therefore
\begin{align}
    S(AB)+S(BC)-S(B)-S(ABC)
        &= \log 2
         > 0.
    \label{eq:cpsv412_strong_subadditivity_defect}
\end{align}
This proves the state-specific Markov obstruction without using the normalized
fixed-point maps.

The full GSNNCH exclusion already concerns the normalized family because
\leanid{MPOTensor.IsGSNNCH} quantifies over
\leanid{MPOTensor.normalizedMPO}. It is now proved directly for the printed
tensor. At four sites, the sector sum in Definition~4.8 of
Ref.~\cite{Cirac2016MPDO_arXiv} splits into positive factors built from the
complementary adjacent-bond products $B_3^{(x)}B_0^{(x)}$ and
$B_1^{(x)}B_2^{(x)}$, acting on the three-site windows $(3,0,1)$ and
$(1,2,3)$, respectively. Orthogonality of distinct sectors removes the cross
terms in their product. Commutativity within each sector identifies the
product of the two factors with the corresponding four-bond sector product.
The resulting positive overlapping factors commute and give the quantum
Markov decomposition
\leanid{MPOTensor.nonempty_quantumMarkovDecomposition_fourCycle_of_isGSNNCHAt},
and hence the strong-subadditivity equality
\leanid{MPOTensor.isSSAEquality_fourCycle_of_isGSNNCHAt}. This derivation uses
only Definition~4.8 of Ref.~\cite{Cirac2016MPDO_arXiv}. It does not assume
zero correlation length, normality, injectivity, primitivity, a single sector,
or completeness of the sector projections.

The identity
\leanid{MPOTensor.CPSVExample412Literal.fourCycleTripartiteState_eq_generic}
shows that this general regrouping is exactly the state used in the literal
entropy calculation. Consequently,
\begin{align}
    \neg\operatorname{IsGSNNCH}(M)
        &\quad\text{holds}.
    \label{eq:cpsv412_literal_not_gsnnch}
\end{align}
This statement is proved by
\leanid{MPOTensor.CPSVExample412Literal.M_not_isGSNNCH}. It discharges the
missing and formerly unlinked four-cycle task tracked in \ghissue{6072} and
the missing Example~4.12 conclusion tracked in \ghissue{6045}. The earlier
attempt to pass through Proposition \emph{prop4to2} had proof-path drift:
that proposition concerns a different condition-level implication and is not
used here. Transport to the separately named $\widehat M$ is unnecessary
because the GSNNCH predicate already uses the normalized family.

\section{Verdict}

The tensor $M$ printed in Example~4.12 of
Ref.~\cite{Cirac2016MPDO_arXiv} is the formal tensor
\leanid{MPOTensor.CPSVExample412Literal.M}. Its positive-length operator
formula in Eq.~\ref{eq:cpsv412_literal_periodic_family} is proved by
\leanid{MPOTensor.CPSVExample412Literal.rho_eq_finKronecker}, and its
saturation statement is proved by
\leanid{MPOTensor.CPSVExample412Literal.M_isSAL}. The transfer calculation in
Eq.~\ref{eq:cpsv412_literal_transfer_scaled_idempotence} is formalized by
\leanid{MPOTensor.CPSVExample412Literal.physTraceTransfer_M_sq}, while
\leanid{MPOTensor.CPSVExample412Literal.physTraceTransfer_M_not_idempotent}
proves Eq.~\ref{eq:cpsv412_literal_transfer_not_idempotent}. Therefore,
\leanid{MPOTensor.CPSVExample412Literal.M_not_isRFPViaTS} proves that the
fixed-representative renormalization assertion is false for the tensor exactly
as printed.

The normalized tensor $\widehat M=(1/2)M$ is the likely local correction. Its
periodic family is identified with the normalized periodic family of $M$ by
tensor-level scalar rescaling and \leanid{MPOTensor.mpo_smul}; it is not an
independent tensor presentation. The parity Kraus families in
Eqs.~\ref{eq:cpsv412_coarse_graining_kraus} and
\ref{eq:cpsv412_refinement_kraus} prove the bare trace-preserving channel
predicate without additional hypotheses, discharging the mathematical task
tracked in \ghissue{6044}. They do not repair the incompatible source
line-246 unit-weight convention. The conclusion therefore remains separated
into the false literal claim, the true normalized channel equations, the
unresolved unit-weight boundary, and the proved full GSNNCH exclusion.

For the normalized four-site state,
\leanid{MPOTensor.CPSVExample412Literal.fourCycleTripartiteState_not_isSSAEquality}
proves the strict defect in
Eq.~\ref{eq:cpsv412_strong_subadditivity_defect}. Its proof uses the exact
marginals
\leanid{MPOTensor.CPSVExample412Literal.traceC_fourCycleTripartiteState},
\leanid{MPOTensor.CPSVExample412Literal.traceA_fourCycleTripartiteState}, and
\leanid{MPOTensor.CPSVExample412Literal.traceAC_fourCycleTripartiteState},
together with their entropy calculations. This discharges the mathematical
task tracked in \ghissue{6073}. The general four-cycle Markov implication and
the exact coordinate identity then give
\leanid{MPOTensor.CPSVExample412Literal.M_not_isGSNNCH}, discharging the tasks
tracked in \ghissue{6072} and \ghissue{6045} without hypotheses beyond the
source GSNNCH definition.

\bibliographystyle{alpha}
\bibliography{references}

\end{document}
